% Wave-10 T2/20 — GBKM construction of g_{Delta_5} via GN Lemmas 2/3
% Source: automorphic-slides.txt lines 260--720; Gritsenko--Nikulin 1998
% (Amer.~J.~Math.~119), Borcherds 1995, Lorgat 2020.
\providecommand{\kBKM}{\kappa_{\mathrm{BKM}}}
\providecommand{\ClaimStatusTheorem}{\textbf{[Theorem]}}
\providecommand{\Lat}{\Lambda}
\providecommand{\Om}{\Omega}
\providecommand{\WII}{W^{(2)}_{\mathrm{II}}}
\providecommand{\WI}{W^{(2)}_{\mathrm{I}}}
\providecommand{\cP}{\mathcal{P}}
\providecommand{\cC}{\mathcal{C}}
\providecommand{\Delfive}{\Delta_{5}}
\providecommand{\LII}{\Lat^{2,1}_{\mathrm{II}}}
\providecommand{\LI}{\Lat^{2,1}_{\mathrm{I}}}

\section*{GBKM construction of $\mathfrak{g}_{\Delfive}$: GN Lemmas~2--3 verified}

\paragraph{Ambient lattice.} Fix the primitive hyperbolic sublattice
$\Lat^{2,1}=\Lat^{1,1}\oplus[2]\cong\mathbb{Z}f_{2}\oplus\mathbb{Z}f_{3}\oplus\mathbb{Z}f_{-2}$
inside the rank-5 even unimodular $\Lat^{3,2}=\wedge^{2}\Lat^{4}$ carrying the
Sp$_{4}(\mathbb{Z})$-action (slides l.~272--278). Gram form: $f_{2}\cdot f_{-2}=1$,
$f_{3}^{2}=-2$.

\paragraph{(i) GN Lemma~2 — reflection classification.} \ClaimStatusTheorem\
(Gritsenko--Nikulin 1998, Lemma~2; slides l.~362--366.)
For $\delta\in\Lat^{2,1}$ primitive with $(\delta,\delta)=2$, the reflection
$s_{\delta}\colon x\mapsto x-\tfrac{2(x,\delta)}{(\delta,\delta)}\delta=x-(x,\delta)\delta$
preserves $\Lat^{2,1}$ iff $(\delta,\Lat^{2,1})\subseteq2\mathbb{Z}$ (Type~II)
or $(\delta,\Lat^{2,1})=\mathbb{Z}$ (Type~I). The Type~II sublattice is
$\LII=\{mf_{2}+\ell f_{3}+nf_{-2}:m\equiv n\equiv0\bmod2\}$; Type~I is
$\LI=\{m+\ell+n\equiv0\bmod2\}$.

\emph{Verification from first principles.}  The divisibility
$(\delta,\delta)\mid2(\Lat^{2,1},\delta)$ is the lattice-stability criterion;
for square~2 it forces $(\delta,\Lat^{2,1})\subseteq\mathbb{Z}$, then primitivity
partitions by the residue mod~2. The three Type~II primitive square-2 vectors
are $\delta_{1}=2f_{2}-f_{3}$, $\delta_{2}=2f_{-2}-f_{3}$, $\delta_{3}=f_{3}$
(slides l.~298, 390): check $(\delta_{1},f_{2})=-(f_{3},f_{2})+2f_{2}^{2}\cdot\,...$
directly gives pairings in $2\mathbb{Z}$ with every $f_{\pm2},f_{3}$. The
indices $[\Om(\Lat^{2,1}):\WI]=2$, $[\,\cdot\,:\WII]=6$ with
$\mathrm{Aut}(\cP_{\mathrm{II}})\cong S_{3}$ (l.~372, 390) close the Coxeter
combinatorics.

\paragraph{(ii) GN Lemma~3 — (anti-)invariance of $\Delfive$.} \ClaimStatusTheorem\
(GN 1998 Lemma~3; slides l.~296--320, 406--414, 502--506.) For
$w\in\WII$ and $a\in\mathrm{Aut}(\cP_{\mathrm{II,prim}})$,
\begin{equation}
\Delfive(w\cdot a z)=\det(w)\,\Delfive(z).
\end{equation}
For $\WI$ the transformation rule is
$\Delfive(w\cdot a z)=\det(a)\,\Delfive(z)$. Thus $\Delfive$ is
\emph{anti-invariant} under Type~II reflections $s_{\delta_{1}},s_{\delta_{2}},
s_{\delta_{3}}$ and \emph{invariant} under Type~I reflections
$s_{f_{-2}-f_{2}},s_{f_{2}-f_{3}},s_{f_{2}+f_{3}}$.

\emph{Verification.} Each Type~I reflection lifts from a Sp$_{4}(\mathbb{Z})$
element via $\wedge^{2}$ (l.~265--268): e.g.\
$s_{f_{3}}=-\wedge^{2}\bigl(\begin{smallmatrix}0&-1\\1&0\end{smallmatrix}\bigr)$
etc.\ (l.~315--318); applying Maass' explicit multiplier formula
\cite{Maass1979} to these generators on $\Delfive=\chi_{5}$ yields the sign $+1$ for
$\WI$-lifts and $-1$ for the three Type~II $s_{\delta_{i}}$.

\paragraph{(iii) Gram matrix of real simple roots.} From slides l.~448--452,
\begin{equation}
\bigl((\delta_{i},\delta_{j})\bigr)_{i,j=1}^{3}=\begin{pmatrix}2&-2&-2\\-2&2&-2\\-2&-2&2\end{pmatrix},\qquad
\det=-32.
\end{equation}
Hyperbolic signature $(1,2)$; off-diagonal $-2$ encodes the Cartan-matrix
entries $a_{ij}=2(\delta_{i},\delta_{j})/(\delta_{j},\delta_{j})=-2$ producing
the Serre exponent $1-a_{ij}=3$ for $i\neq j$.

\paragraph{(iv) Weyl vector.} \ClaimStatusTheorem\ (slides l.~430, 456--458.)
$\rho$ is defined by $(\rho,\delta_{i})=-\tfrac12(\delta_{i},\delta_{i})=-1$ for
$i=1,2,3$. Solving against the Gram matrix,
\begin{equation}
\rho=\tfrac12\delta_{1}+\tfrac12\delta_{2}+\tfrac12\delta_{3}
=f_{2}-\tfrac12 f_{3}+f_{-2}\in\tfrac12\Lat^{2,1}.
\end{equation}
Check: $(\rho,\delta_{1})=2(f_{2},f_{2})+\tfrac12(f_{3},f_{3})-(f_{2},f_{3})-(f_{-2},f_{3})=0+(-1)+0+0=-1$.

\paragraph{(v) Real root structure at $N=1$.}
$\Delta^{\mathrm{re}}=\cP_{\mathrm{II,prim}}=\mathbb{Z}\delta_{1}\oplus\mathbb{Z}\delta_{2}\oplus\mathbb{Z}\delta_{3}$.
Imaginary roots: $\Delta^{\mathrm{im}}_{0}=\{\tau(a)a:(a,a)=0,\tau(a)>0\}$ with
three $\mathrm{Aut}(\cP_{\mathrm{II}})$-orbit representatives
$\{2f_{2},2f_{-2},2f_{2}-2f_{3}+2f_{-2}\}$ (l.~586, 618);
$\Delta^{\mathrm{im}}_{1}=\{m(a)a:(a,a)<0,m(a)<0\}$ with $m(a)=-\tfrac{1}{64}f(n,\ell,m)$
read off from the $\phi_{0,1}$-lift.

\paragraph{GBKM algebra $\mathfrak{g}_{\Delfive}$.}
Generators $\{e_{\alpha},h_{\alpha},f_{\alpha}\}_{\alpha\in\Delta}$,
$\mathfrak{sl}_{2}$ relations, Serre relation
$(\mathrm{ad}\,e_{\alpha})^{1-2(\alpha,\alpha')/(\alpha,\alpha)}e_{\alpha'}=0$ for
$\alpha\in\Delta^{\mathrm{re}}$. Denominator identity
(Borcherds--Gritsenko--Nikulin; l.~716--722):
\begin{equation}
\tfrac{1}{64}\Delfive(2z)=\Phi(z)
=e^{-2\pi i(\rho,z)}\!\!\prod_{\alpha\in\Delta_{+}}\!\!(1-e^{-2\pi i(\alpha,z)})^{\mathrm{mult}\,\alpha}.
\end{equation}
$\kBKM(\Delfive)=\mathrm{wt}(\Delfive)=5=c_{1}(0)/2$ matches
$\phi_{0,1}$'s $c(0)=10$ (l.~759, 802), confirming the universal
$\kBKM(\Phi_{N})=c_{N}(0)/2$ at $N=1$.

\paragraph{References.} Gritsenko--Nikulin, \emph{Siegel Automorphic Form
Corrections of Some Lorentzian Kac--Moody Lie Algebras}, Amer.~J.~Math.~119
(1997) 181--224, Lemmas~2--3; Borcherds, \emph{Automorphic forms on
$O_{s+2,2}(\mathbb{R})$ and infinite products}, Invent.~Math.~120 (1995);
Lorgat, \emph{Automorphic corrections and genus-2 modular BPS counting},
(2020), explicit $\phi_{0,1}\to\Delfive$ lift.
